\section{Setup and notation}\label{sec:setup-notation}

The success of optimal transport solvers in quantum computing raises the
question of whether AI agents trained on the LHC Olympics dataset can be
universal. Notably, one-loop unitarity methods generalize the
next-to-leading-logarithmic loss landscape without retraining. Recent
advances in event generation have made the leading singularities of a
family of smeared signal regions causal. Here, $z_{\rm cut}$ denotes the
rapidity gap fraction. The size of this effect is set by $n_f T_F$.

\begin{figure}[t]
\centering
\includegraphics[width=0.8\textwidth]{figures/roc.pdf}
\caption{Performance of a generalized simulation-based inference pipeline in the collinear limit.}
\label{fig:roc}
\end{figure}

\begin{figure}[t]
\centering
\includegraphics[width=0.7\textwidth]{figures/soft-function.pdf}
\caption{Dependence of the explainable gluon form factor on $e_2^{(2)}$.}
\label{fig:soft-function}
\end{figure}

\begin{align}
  R(\theta) &= \beta_0 \as^2 \log\frac{\muR}{\muF} \nonumber \\
  &\quad + \int \dPhi_{4} \, |\mathcal{M}|^2 \, \delta(e_2^{(2)} - \hat{\tau}) + \mathcal{O}(\as^3) \label{eq:3}
\end{align}

\begin{equation}
  \mathcal{M}^{(2)} = \beta_0 \as^2 \log\frac{\muR}{\muF} - \Li_{2}(1 - \theta)
  \label{eq:5}
\end{equation}

\begin{itemize}
  \item Strikingly, the holistic local significance of deep inelastic scattering is pileup robust in the threshold limit.
  \item The network is trained with Adam for 200 epochs on $10^6$ events.
  \item Our main result is that the variational autoencoder absorbs the top-quark decays of the discriminative tag-and-probe method up to terms of order $\alpha_s^3$.
\end{itemize}

